\documentclass[10pt]{article}
\providecommand{\FIGDIR}{fig}
\newcommand{\DIGESTDATE}{1 September 2026}
% ===================== PM Research Watch - shared preamble =====================
\usepackage[T1]{fontenc}
\usepackage[utf8]{inputenc}
\usepackage{amsmath,amssymb}
\usepackage{tgpagella}
\usepackage{tgheros}
\renewcommand{\sfdefault}{qhv}
\renewcommand{\ttdefault}{lmtt}
\usepackage{microtype}
\usepackage[english]{babel}

\usepackage[a4paper,top=17mm,bottom=17mm,left=17mm,right=17mm,
            headsep=6mm,footskip=10mm]{geometry}
\usepackage{graphicx}
\usepackage{xcolor}
\usepackage{colortbl}
\usepackage{tikz}
\usetikzlibrary{calc}
\usepackage[most]{tcolorbox}
\usepackage{titlesec}
\usepackage{fancyhdr}
\usepackage{booktabs}
\usepackage{array}
\usepackage{longtable}
\usepackage{enumitem}
\usepackage{multicol}
\usepackage{ragged2e}
\usepackage[hidelinks]{hyperref}

% ------------------------------------------------------------------ palette
\definecolor{Ink}    {HTML}{14262E}
\definecolor{Deep}   {HTML}{0F4C5C}
\definecolor{Teal}   {HTML}{1F7A8C}
\definecolor{Sky}    {HTML}{4FA3C4}
\definecolor{Sage}   {HTML}{6FA287}
\definecolor{Amber}  {HTML}{E3A008}
\definecolor{Clay}   {HTML}{B5643C}
\definecolor{Coral}  {HTML}{D1495B}
\definecolor{Violet} {HTML}{6A5B9E}
\definecolor{Slate}  {HTML}{4A6572}
\definecolor{Paper}  {HTML}{FBF9F5}
\definecolor{Mist}   {HTML}{EEEAE2}
\definecolor{Rule}   {HTML}{D9D3C9}

\pagecolor{Paper}
\color{Ink}

\hypersetup{colorlinks=true, linkcolor=Deep, urlcolor=Teal, citecolor=Teal,
            pdftitle={PM Research Watch}, pdfauthor={Automated daily digest}}

% ------------------------------------------------------------------ headings
\titleformat{\section}
  {\sffamily\bfseries\large\color{Deep}}{}{0pt}
  {}[{\vspace{1.2mm}\color{Amber}\titlerule[1.1pt]}]
\titlespacing*{\section}{0pt}{5.5mm}{2.8mm}

\titleformat{\subsection}
  {\sffamily\bfseries\normalsize\color{Teal}}{}{0pt}{}
\titlespacing*{\subsection}{0pt}{3.6mm}{1.4mm}

\setlist[itemize]{leftmargin=4.2mm, itemsep=1.1mm, topsep=1.3mm, parsep=0pt,
                  label=\textcolor{Teal}{\footnotesize$\blacktriangleright$}}

% ------------------------------------------------------------------ header/footer
\pagestyle{fancy}
\fancyhf{}
\renewcommand{\headrulewidth}{0pt}
\renewcommand{\footrulewidth}{0pt}
\fancyhead[L]{\sffamily\scriptsize\color{Slate}PM RESEARCH WATCH\;\textcolor{Rule}{\textbar}\;\DIGESTDATE}
\fancyhead[R]{\sffamily\scriptsize\color{Slate}Particulate matter monitoring \& health effects}
\fancyfoot[C]{\sffamily\scriptsize\color{Slate}\thepage}

% ------------------------------------------------------------------ components
% metric tile
\newcommand{\metric}[3]{%
  \begin{tcolorbox}[enhanced, colback=white, colframe=Rule, boxrule=0.5pt,
    arc=1.6mm, left=2.2mm, right=2.2mm, top=1.8mm, bottom=1.8mm,
    borderline west={1.7mm}{0pt}{#3}, nobeforeafter, valign=center,
    height=17mm, width=\linewidth]
    \RaggedRight
    {\sffamily\bfseries\fontsize{17}{17}\selectfont\color{#3}#1}\\[0.5mm]
    {\sffamily\fontsize{6.5}{7.4}\selectfont\color{Slate}#2}
  \end{tcolorbox}}

% pull-quote / key insight box
\newtcolorbox{keybox}[1][]{enhanced, colback=Mist, colframe=Mist,
  arc=1.6mm, left=3.2mm, right=3.2mm, top=2.4mm, bottom=2.4mm,
  borderline west={1.7mm}{0pt}{Amber}, #1}

% subtopic band
\newcommand{\band}[2]{%
  \vspace{3.2mm}\noindent
  \begin{tcolorbox}[enhanced, colback=#2, colframe=#2, arc=1.2mm,
    left=2.6mm, right=2.6mm, top=1.1mm, bottom=1.1mm, boxrule=0pt,
    width=\linewidth, nobeforeafter]
    {\sffamily\bfseries\small\color{white}#1}
  \end{tcolorbox}\vspace{1.6mm}}

% one paper entry: {short cite}{title}{journal}{doi}{body}
\newcommand{\paperentry}[5]{%
  \par\vspace{1.9mm}\noindent
  {\sffamily\bfseries\footnotesize\color{Ink}#1}\;%
  {\sffamily\fontsize{7}{8}\selectfont\color{Slate}#3}\\[0.4mm]
  {\footnotesize\itshape\color{Deep}#2}\\[0.6mm]
  {\footnotesize #4}\\[0.3mm]
  {\fontsize{6.6}{7.6}\selectfont\color{Slate}\texttt{doi:}\,\href{https://doi.org/#5}{\texttt{\detokenize{#5}}}}
  \par}

% caption style
\newcommand{\figcap}[1]{\par\vspace{1.0mm}
  {\fontsize{7.2}{8.6}\selectfont\color{Slate}\textsf{\textbf{Fig.}}\; #1\par}\vspace{1.5mm}}

\providecommand{\permil}{\textperthousand}

\setlength{\parindent}{0pt}
\setlength{\parskip}{0pt}
\renewcommand{\arraystretch}{1.18}


\begin{document}

% ==================================================================== MASTHEAD
\thispagestyle{empty}
\begin{tikzpicture}[remember picture, overlay]
  \fill[Deep] (current page.north west) rectangle
      ($(current page.north east)+(0,-67mm)$);
  \fill[Amber] ($(current page.north west)+(0,-67mm)$) rectangle
      ($(current page.north east)+(0,-68.4mm)$);
  \foreach \x/\y/\r in {22/12/0.9, 41/26/1.5, 63/9/0.7, 88/31/1.1, 112/17/1.9,
                        138/29/0.8, 159/13/1.3, 178/54/1.0, 196/21/1.6,
                        30/44/0.6, 74/57/1.2, 125/48/0.9, 168/6/0.8, 191/59/0.7,
                        52/18/1.0, 100/60/0.6, 148/20/1.1, 66/49/1.4, 205/45/0.8}
    \fill[Sky, opacity=0.30] ($(current page.north west)+(\x mm,-\y mm)$) circle (\r mm);
\end{tikzpicture}

\vspace*{4mm}
{\sffamily\fontsize{7.6}{9}\selectfont\color{Sky}%
 AUTOMATED DAILY LITERATURE DIGEST\;\textbullet\;PhD RESEARCH WATCH\par}
\vspace{2.2mm}
{\sffamily\bfseries\fontsize{27}{29}\selectfont\color{white}%
 Particulate Matter\par}
{\sffamily\fontsize{27}{29}\selectfont\color{Sky}%
 Monitoring \& Health Effects\par}
\vspace{5.0mm}
{\sffamily\fontsize{8.4}{10}\selectfont\color{white}%
 Issue 2026-09-01\;\;\textcolor{Amber}{\textbar}\;\;Entry window 1 September 2026 (single day)\;\;%
 \textcolor{Amber}{\textbar}\;\;33 records in scope, 98 screened out\par}
\vspace{18mm}

% -------------------------------------------------------------- metric strip
\noindent
\begin{minipage}[t]{0.190\linewidth}\metric{33}{RECORDS IN SCOPE\\(98 SCREENED OUT)}{Deep}\end{minipage}\hfill
\begin{minipage}[t]{0.190\linewidth}\metric{9}{SUBTOPIC CLUSTERS\\OF TEN OCCUPIED}{Teal}\end{minipage}\hfill
\begin{minipage}[t]{0.190\linewidth}\metric{16}{EXPOSURE ASSESS.\\\& MODELLING}{Amber}\end{minipage}\hfill
\begin{minipage}[t]{0.190\linewidth}\metric{1}{EFFECT ESTIMATE\\WITH A RATIO CI}{Clay}\end{minipage}\hfill
\begin{minipage}[t]{0.190\linewidth}\metric{22}{RECORDS WITH NO\\HEALTH ENDPOINT}{Sky}\end{minipage}

\vspace{6mm}

% -------------------------------------------------------------- signal of day
\section*{\sffamily\bfseries\color{Deep}Signal of the day}
\vspace{-1.5mm}{\color{Amber}\rule{\linewidth}{1.1pt}}\vspace{2.5mm}

\begin{keybox}
{\sffamily\bfseries\small\color{Deep}Two independent \textit{Atmosphere} papers, published
the same day and citing nothing in common, make the same argument: the averaging interval
written into air-quality regulation destroys most of what a high-rate network measures.}\\[1.4mm]
\footnotesize
Raso et al.\ ran certified reference instruments alongside calibrated low-cost sensors
through two consecutive New Year's Eve firework events. Conventional monitoring detects
that an episode happened; it does not characterise it. Minute-resolution PM$_{10}$ reached
\textbf{215\,$\mu$g\,m$^{-3}$}, roughly three times the corresponding hourly value ---
a peak amplification factor of \textbf{2.90--2.93} that the hourly channel simply does not
contain. Mohd Nadzir et al.\ deploy the complementary case: a single low-cost multi-sensor
node in the Antarctic Peninsula, \textbf{58{,}652 valid minute observations} over
May--June 2023, in an environment where the daily mean is near-constant and therefore
uninformative. At minute resolution the record resolves \textbf{142 single-pollutant and
30 multi-pollutant enhancement events}; the median PM$_{2.5}$/PM$_{10}$ ratio of
\textbf{0.96} says the mass is essentially all fine, which is the diagnostic that
distinguishes camp and vessel combustion from sea-salt and local dust.
Both frame the same brief from opposite ends of the concentration range ---
215\,$\mu$g\,m$^{-3}$ urban peaks against a 6.3\,$\mu$g\,m$^{-3}$ polar background ---
and both land on temporal resolution, not accuracy, as the property that justifies the
sensor. That is the strongest argument this watch has carried for the distributed low-cost
network as an instrument in its own right rather than a cheap approximation to a reference
station. Neither reports a drift or co-location series long enough to say what those
minute values are worth after six months in the field.
\end{keybox}

\vspace{3mm}

% -------------------------------------------------------------- exec brief
\section*{\sffamily\bfseries\color{Deep}What changed today}
\vspace{-1.5mm}{\color{Amber}\rule{\linewidth}{1.1pt}}\vspace{2.5mm}

\footnotesize
\begin{itemize}
\item \textbf{The largest corpus since 27 August, and the most lopsided.}
\textit{Exposure assessment \& modelling} takes \textbf{16 of 33} records and
\textbf{22 of 33 (67\%) carry no human health endpoint at all}. Seven of the sixteen
arrive in one \textit{J.\ Environ.\ Sci.} deposit batch, so the cluster's size is partly a
publishing artefact, not a shift in the field.

\item \textbf{A policy evaluation with an individual-level counterfactual.} Sun D et al.\
(\textit{Lancet Planet Health}) use China Kadoorie Biobank participants ($n$=34{,}862)
assigned to intervention or control by their municipality's PM-reduction target under the
2013 Clean Air Act. The intervention group's predicted 10-year CVD risk rose
\textbf{3.95\% (95\% CI 3.18--4.72) less} than control; per 10\,$\mu$g\,m$^{-3}$,
\textbf{1.80 (1.34--2.27) percentage points} for PM$_{2.5}$. The outcome is a
\emph{predicted} risk score, not incident disease, and assignment is a municipal policy
target --- but this is the closest thing to a quasi-experiment the PM--CVD literature has
at this scale.

\item \textbf{Only one estimate in the whole corpus is a ratio with a confidence
interval.} G\'omez-Olarte et al.\ report \textbf{aOR 1.50 (1.02--2.23)} for bronchitis at
age one in boys per IQR (2.0\,$\mu$g\,m$^{-3}$) PM$_{2.5}$, with IL-8 up
\textbf{31.9\% (18.6--46.7)} in girls. Today's forest plot is therefore a single point,
not a synthesis; the Lancet numbers are percentage-point differences on a risk score and
cannot share a log-ratio axis with it.

\item \textbf{Two records converge on the same conclusion about PM mass as a metric.}
Arora et al.'s commentary restates the source-dependence of oxidative potential ---
biomass-burning particles up to \textbf{35$\times$} the OP of suburban PM$_{2.5}$, O$_3$
aging generally \emph{reducing} it --- while Arpornwichanop \& Chuersuwan reach the same
place from the elemental side: deposited intake in Nonthaburi ranks Al, Fe, B, Zn highest,
the solubility- and toxicity-adjusted priority ranks \textbf{As, Cd, Sb}.

\item \textbf{Fine-scale exposure modelling posts a concrete number.} Wang T et al.'s
IAQMS-Industry resolves PM$_{2.5}$ at \textbf{100\,m} inside industrial parks and
attributes \textbf{22.9--26.4\%} of park PM$_{2.5}$ to local point sources against
1.6--13.7\% for the same emissions in a 3\,km grid, with normalised mean bias improving
from $-$16.9\%\,--\,$-$7.7\% to \textbf{3.1--6.2\%}. That gap is the sub-grid smoothing
error a dense network would be measuring against.

\item \textbf{Sub-Saharan Africa and Latin America are both absent again}; China supplies
13 of 33, Europe 6, North America 2. \textbf{Source breakages:} OpenAlex HTTP 403
(plan-gated date filter, unchanged since 28 July) and arXiv HTTP 429 on four attempts, so
the preprint-instrumentation leg contributed nothing. All three Consensus hits on the
sensor-calibration axis were already in the store (2\,Aug, 5\,Aug, 30\,Jul).
\end{itemize}

\clearpage

% ==================================================================== ANALYTICS
\section*{\sffamily\bfseries\color{Deep}Corpus analytics}
\vspace{-1.5mm}{\color{Amber}\rule{\linewidth}{1.1pt}}\vspace{3mm}

\noindent\includegraphics[width=\linewidth]{\FIGDIR/f1_subtopics.png}
\figcap{Subtopic assignment of the 33 in-scope records. Nine of ten clusters are occupied;
only \textit{Other clinical endpoints} is empty. \textit{Exposure assessment \& modelling}
at 16 is the largest single-cluster count in any single-day issue since the 26 July seed
issue (26, same cluster), and the seven \textit{J.\ Environ.\ Sci.} papers inside it
(soot morphology, brown-carbon triplets, aqSOA, NH$_3$--isoprene, terrain blocking,
N-containing organics, industrial-park CTM) are one publisher batch rather than seven
independent signals.}

\vspace{2mm}
\noindent\includegraphics[width=\linewidth]{\FIGDIR/f2_design_endpoint.png}
\figcap{Left --- study architecture. \textit{Modelling / inventory} (12) and
\textit{Measurement campaign} (8) together hold 20 of 33; four cohorts, five
laboratory/toxicology studies, one cross-sectional, one intervention and one review make
up the rest. Right --- health endpoint. \textbf{Twenty-two of thirty-three records (67\%)
carry no human health endpoint} --- level with 9 August and behind only 20 August (72\%),
2 August (71\%) and 5 August (70\%). Of the eleven that do, four are cancer-risk screens
rather than observed incidence.}

\vspace{2mm}
\noindent\includegraphics[width=\linewidth]{\FIGDIR/f3_metric_geography.png}
\figcap{Left --- particle metric. PM$_{2.5}$ alone is the exposure in 8 records and
appears jointly with PM$_{10}$ in 6 more; three records are optical (soot morphology,
Mediterranean dust, N-containing organic absorbance) and one is bioaerosol. The 13 in
\textit{PM (unspeciated) / emissions} are mostly the chemistry and inventory papers.
Right --- geographic provenance. China at 13 is 39\% of the corpus. \textbf{Sub-Saharan
Africa and Latin America are both absent}, as they were on 31 August; the two
\textit{Polar / remote marine} entries are the Antarctic sensor node and the CAFE-Pacific
aircraft campaign.}

\vspace{2mm}
\noindent\includegraphics[width=\linewidth]{\FIGDIR/f6_lifecourse.png}
\figcap{Life-course windows addressed; counts are non-exclusive and only the six records
studying identifiable human populations contribute. The \textbf{65+ cell is zero}: no
record in this issue reports an age-stratified estimate for older adults, even though the
Kadoorie cohort (mean baseline age 51.3, followed to 2020--21) certainly contains a large
older stratum. Scoring it otherwise would misdescribe what the paper reports.}

\vspace{2mm}
\noindent\includegraphics[width=\linewidth]{\FIGDIR/f4_heatmap.png}
\figcap{Subtopic $\times$ study-architecture cross-tabulation. The separation this watch
keeps flagging is at its widest today: \textit{Exposure assessment \& modelling} and
\textit{Sensing, forecasting \& instrumentation} together hold \textbf{21 of 33 records
and contribute exactly zero effect estimates}. The four cohort records span only three
health clusters --- two of them the same one (\textit{Neuro / mental health}) --- and the
single cardiovascular record is not a cohort but a difference-in-differences design.}

\vspace{2mm}
\noindent\includegraphics[width=\linewidth]{\FIGDIR/f5_forest.png}
\figcap{\textbf{A single estimate.} Only G\'omez-Olarte et al.\ report an effect on a
ratio scale with a confidence interval, so this panel displays one contrast and is not a
synthesis. The lower bound (1.02) sits against the null, the cohort is small ($n$=606 for
the respiratory outcome), and the estimate is sex-stratified post hoc --- read it as
consistent with the cytokine gradient in the same paper, not as independent confirmation
of it.}

\clearpage

% ==================================================================== DIGEST
\section*{\sffamily\bfseries\color{Deep}Paper-level digest}
\vspace{-1.5mm}{\color{Amber}\rule{\linewidth}{1.1pt}}\vspace{1mm}

{\fontsize{7.4}{9}\selectfont\color{Slate}Ordered by cluster. Each entry gives the
methodological core and the finding that matters, not an abstract paraphrase.\par}

\band{Neuro / mental health\hfill 2 records}{Deep}

\paperentry{Barbalat et al.\ 2026}
{Heat, cold and ambient pollution in the first 1000 days and autism risk indicators}
{Environ Sci Technol}
{ELFE national birth cohort, $n$=12{,}139 two-year-olds, M-CHAT-R screening score against
weekly temperature and PM$_{2.5}$/PM$_{10}$/NO$_2$ from 1\,km daily predictions (200\,m
for NO$_2$ and temperature in large urban areas), modelled with confounder-adjusted
distributed lag non-linear models. \textbf{Prenatal air pollution was not associated with
the screening score}; night-time heat in the early third trimester was, and survived
adjustment for PM$_{2.5}$ and PM$_{10}$. Postnatally the pattern is sex-specific: higher
PM$_{2.5}$ and PM$_{10}$ at months 7--11.5 raised scores in males relative to females.
The outcome is a screening instrument, not a diagnosis, and the sex-stratified postnatal
windows are the kind of finding a DLNM will produce by construction unless
pre-registered. Useful as a negative result on the prenatal PM window.}
{10.1021/acs.est.5c18062}

\paperentry{Kuznetsov et al.\ 2026}
{Particulate matter and life satisfaction, across and within twin pairs}
{Environ Res}
{TwinLife, $n$=4{,}504 young-adult twins (2{,}249 MZ / 2{,}255 DZ), mean age 23.0, up to
six waves 2014--2022, decomposing the PM--wellbeing association into between- and
within-pair components. The effect runs \textbf{between families} ($b$=$-$0.05,
SE 0.01, $p<$.001) and strengthens with age and calendar year for PM$_{10}$. Within
pairs, after removing shared family variance, the association is null in DZ twins and
only nominally significant in MZ twins for PM$_{10}$ ($b$=$-$0.09, SE 0.04, $p$=0.04)
and not at all for PM$_{2.5}$ ($p$=.15). That is the honest reading: most of the apparent
PM effect on subjective wellbeing is confounded by family SES and neighbourhood, and the
design is what makes that visible. A rare case of a study whose null is the contribution.}
{10.1016/j.envres.2026.125565}

\band{Cardiovascular \& metabolic\hfill 1 record}{Teal}

\paperentry{Sun D et al.\ 2026}
{China's Clean Air Act and individual-level cardiovascular risk}
{Lancet Planet Health}
{China Kadoorie Biobank, $n$=34{,}862 CVD-free at baseline, mean age 51.3, split into
intervention ($n$=25{,}497) and control ($n$=9{,}365) by the local government's PM
reduction target under the 2013 Clean Air Act, with 10-year predicted CVD risk as the
outcome and a difference-in-differences estimator. The intervention group's predicted
risk rose \textbf{3.95\% (3.18--4.72) less} than control, with larger differences under
stricter enforcement. Per 10\,$\mu$g\,m$^{-3}$, 2013--2021: PM$_{2.5}$ \textbf{1.80
(1.34--2.27)} percentage points, PM$_{10}$ 1.24 (0.84--1.63), O$_3$ 0.58 (0.33--0.83);
overall pollutant change corresponded to a 6.6 percentage-point reduction in predicted
risk. Two real limits: the outcome is a validated \emph{prediction}, not observed
incidence, and treatment assignment is a municipal policy target, so residual regional
confounding is structural rather than adjustable. O$_3$ rose in both arms.}
{10.1016/j.lanplh.2026.101513}

\band{Reproductive \& developmental\hfill 1 record}{Sky}

\paperentry{Gao et al.\ 2026}
{Urinary biomarkers of airborne organics in pregnancy and small-for-gestational-age risk}
{Environ Sci Technol}
{TIMFEM prospective cohort, $n$=1{,}425 mother--infant pairs, 20 urinary biomarkers of
airborne organic pollutants by LC-MS/MS (p-phenylenediamines, nitro-PAHs,
phenylguanidines, benzothiazoles, cotinine), 18 detected in $>$40\% of participants, with
leave-one-out in vitro screening in human trophoblast cells at exposure-guided doses.
Each ln-unit increase in \textbf{6PPD was associated with 27\% higher SGA risk and 6PPD-Q
with 23\%}; mixture models added TPG, cotinine and 1-aminopyrene. The mechanism the cell
work supports is metabolic rather than cytotoxic --- reduced pyruvate entry into the TCA
cycle and impaired isocitrate-to-$\alpha$-ketoglutarate conversion at doses that barely
touched proliferation. No confidence intervals are given for the epidemiological
estimates in the abstract, so they are not plotted here. Relevance: these are
tyre-and-road and traffic-associated organics that no PM mass channel resolves.}
{10.1021/acs.est.6c04451}

\band{Respiratory \& allergic\hfill 1 record}{Sage}

\paperentry{G\'omez-Olarte et al.\ 2026}
{Early-life air pollution, respiratory outcomes and stimulated cytokine response}
{Environ Pollut}
{German birth cohort, $n$=606 one-year-olds for bronchitis and wheezing, cytokine
subcohort $n$=508 (LPS-stimulated TNF-$\alpha$, MCP-1, IL-8, IL-6, IL-10), residential
PM$_{2.5}$/PM$_{10}$/NO$_2$ per IQR of 2.0/2.8/9.1\,$\mu$g\,m$^{-3}$, with BKMR for the
mixture and stratification by sex and temperature. In boys, PM$_{2.5}$ gave
\textbf{aOR 1.50 (1.02--2.23)} for bronchitis and \textbf{+23.8\% (11.2--37.9)} IL-8;
in girls, \textbf{+31.9\% (18.6--46.7)} IL-8, +12.8\% (3.9--22.6) TNF-$\alpha$ and
+35.7\% (14.4--61.1) IL-10. The mixture association with IL-8 and TNF-$\alpha$ was
non-linear and appeared only above the 60th percentile of all pollutants. Small cohort,
and the sex split doubles the tests; the consistent direction across five cytokines is
what carries it. PM$_{2.5}$, not NO$_2$, is the primary predictor.}
{10.1016/j.envpol.2026.129079}

\clearpage

\band{Mechanistic toxicology\hfill 2 records}{Violet}

\paperentry{Meldrum \& Leonard 2026}
{Diesel particle priming reprograms the airway response to a subsequent allergen}
{Front Immunol}
{Murine model with deliberate temporal separation between diesel exhaust particle
exposure and house dust mite challenge --- so neither co-exposure nor prior
sensitisation --- read out with compartment-resolved bulk transcriptomics (lung tissue
vs airway lumen) plus targeted single-cell profiling. DEP priming amplified leukocyte
recruitment on subsequent HDM challenge and did so \textbf{neutrophilically, with no
eosinophil response}. DEP alone drove NF-$\kappa$B programmes; HDM drove NF-$\kappa$B
and interferon; the primed response was qualitatively different from either, with
elevated ISG expression and expanded neutrophil subpopulations concentrated in the
luminal compartment. The design is the contribution: it isolates a memory-like effect of
a particle exposure that has already cleared. Murine, single particle type, no
dose--response.}
{10.3389/fimmu.2026.1869644}

\paperentry{Arora et al.\ 2026}
{Oxidative potential of fresh versus ozone-aged PM$_{2.5}$ across Chinese sources}
{J Environ Sci (China)}
{A commentary, not primary data --- it summarises Ma et al.\ (\texttt{10.1016/j.jes.2024.04.023}),
which measured DTT-assay oxidative potential on fresh and O$_3$-aged
PM$_{2.5}$ from multiple source types. \textbf{Biomass-burning particles carried up to
35$\times$ the OP of suburban PM$_{2.5}$}, driven by water-soluble organics and
transition metals, and O$_3$ aging generally \emph{lowered} OP while producing complex
chemical transformation. Counted and listed here because the source-dependence claim is
the one this watch tracks, but no number in it is new and the aging direction runs
opposite to the usual assumption that atmospheric processing increases toxicity.
Treat as a pointer to the 2024 primary paper.}
{10.1016/j.jes.2026.05.042}

\band{Occupational \& indoor\hfill 2 records}{Clay}

\paperentry{Yang et al.\ 2026}
{Ventilation optimisation for a 300\,m welding workshop with distributed dust sources}
{Atmosphere}
{CFD of airflow and welding-fume transport in a 300 $\times$ 28 $\times$ 21\,m hall,
varying exhaust-to-supply ratio and side-exhaust-port height, and contrasting low-density
aluminium-alloy fume with high-density carbon-steel fume. Two air-velocity peaks appear
at 0--2\,m and 8--12\,m above floor; the \textbf{top exhaust port outperforms the side
port} throughout, higher ESR accelerates upward migration and suppresses lateral
dispersion, and a badly set side-port height produces either short-circuiting or excess
lateral spread. Optimal side-port height \textbf{5\,m for aluminium, 6\,m for carbon
steel}. Simulation only --- no measured personal exposure, and fume is treated as a
passive density-differentiated tracer rather than a coagulating aerosol.}
{10.3390/atmos17090843}

\paperentry{Jaiswal et al.\ 2026}
{Nature-based indoor filtration in an EDM gear-manufacturing shop}
{Research Square (preprint)}
{Single-site five-month before/after intervention in Gurugram, India. Baseline was
severe: \textbf{PM$_{2.5}$ 330\,$\mu$g\,m$^{-3}$ and CO 88\,ppm}. Plant-based filter
units plus a redesigned vent/exhaust layout, sited using airflow visualisation, cut
\textbf{CO by 90\% and PM$_{2.5}$ by 42\% on average} --- leaving PM$_{2.5}$ around
190\,$\mu$g\,m$^{-3}$, which the authors say plainly is still not a healthy range. No
control period, no controlled comparison against a conventional HEPA unit, and the CO
and PM reductions differ by a factor the paper does not explain mechanistically.
Preprint, not peer reviewed. Listed for the occupational baseline number, which is the
part that is hard to obtain.}
{10.21203/rs.3.rs-10511730/v1}

\band{Sensing, forecasting \& instrumentation\hfill 5 records}{Amber}

\paperentry{Raso et al.\ 2026}
{Temporal resolution and the characterisation of episodic pollution events}
{Atmosphere}
{Two consecutive New Year's Eve firework campaigns with certified reference
instrumentation co-deployed with calibrated low-cost sensors, comparing daily, hourly and
minute aggregation of PM$_{10}$, NO$_2$, CO and O$_3$. Regulatory averaging detects the
episode and misses its structure: minute PM$_{10}$ reached \textbf{215\,$\mu$g\,m$^{-3}$},
a \textbf{peak amplification factor of 2.90--2.93} over the hourly value. The direct
implication for distributed sensing is that the case for a dense network is not
redundancy against reference-grade accuracy but access to a time scale the reference
framework discards --- and that acute-exposure metrics built on hourly data are
systematically low for episodic sources. Two events at one site; the PAF is a description,
not a transferable constant.}
{10.3390/atmos17090860}

\paperentry{Mohd Nadzir et al.\ 2026}
{Minute-scale trace gas and PM variability in the Antarctic Peninsula from a low-cost node}
{Atmosphere}
{A single low-cost multi-sensor package, May--June 2023, \textbf{58{,}652 valid
observations} of NO$_2$, O$_3$, CO, SO$_2$, PM$_{2.5}$ and PM$_{10}$ (TVOC in May only).
Medians: PM$_{2.5}$ \textbf{6.29} and PM$_{10}$ \textbf{7.29\,$\mu$g\,m$^{-3}$},
O$_3$ 38.17\,ppb, NO$_2$ 11.05\,ppb, SO$_2$ 0.01\,ppb. The
\textbf{PM$_{2.5}$/PM$_{10}$ ratio of 0.96} indicates an almost entirely fine-mode
aerosol, consistent with combustion rather than sea salt or crustal dust. Percentile-based
event detection found \textbf{142 single- and 30 multi-pollutant enhancements}. The TVOC
channel was 98\% zeros and was rightly demoted. What is missing is any co-location or
drift assessment --- at these concentrations a few $\mu$g\,m$^{-3}$ of baseline offset
would change every conclusion, and an uncalibrated nephelometric PM channel in high
humidity is exactly where that offset lives.}
{10.3390/atmos17090855}

\paperentry{Alraddadi et al.\ 2026}
{Explainable one-hour-ahead AQI forecasting across eight Saudi cities}
{Atmosphere}
{Attention-based bidirectional recurrent networks against a persistence baseline in eight
climatically distinct cities. The deep models \textbf{matched persistence on RMSE and
$R^2$} (cross-city mean $R^2\approx$0.90) and beat it on MAE by roughly a fifth, with the
margin largest in the most variable cities. Event-based evaluation flagged
\textbf{84--91\% of high-AQI exceedances}. Surface meteorology added little on average and
its value was city- and architecture-specific. SHAP and LIME put particulate matter as the
dominant driver everywhere. The persistence comparison is the part worth copying: at
one-hour horizon hourly AQI is so autocorrelated that an unbenchmarked $R^2$ of 0.90 means
nothing, and most papers in this genre do not report the baseline.}
{10.3390/atmos17090858}

\paperentry{Chen C et al.\ 2026}
{Amine-containing particles from fireworks by single-particle mass spectrometry}
{J Hazard Mater}
{2023 Spring Festival field campaign with single-particle aerosol mass spectrometry plus
source samples. Amine-containing particles rose \textbf{96.2\%} during the firework
period, peak size shifted \textbf{0.42\,$\to$\,0.44\,$\mu$m}, and internal mixing with
sulfate increased; fireworks raised metal and elemental-organic-carbon classes while
depressing biomass-burning and secondary classes. \textbf{N-methylaniline} is proposed as
a tracer separating fireworks from firecrackers. A multi-model comparison identifies
modified aerosol acidity, not liquid water content, as the dominant control on
gas--particle partitioning of amines, with ALWC showing a non-linear threshold. One
festival, one city; the tracer claim needs a second site before it is usable.}
{10.1016/j.jhazmat.2026.143446}

\paperentry{Hong et al.\ 2026}
{Predicting fungal and bacterial bioaerosol concentrations across building types}
{Indoor Air}
{\textbf{9{,}377 indoor air samples} across 27 multiuse facility types in South Korea,
modelled with a generalised additive mixed model to capture non-linear temperature and
humidity effects, then a random forest on the resulting sensitivity classes. Fungi
respond mainly to relative humidity (16 facility types), bacteria mainly to temperature
(7); peak microbial levels occur at \textbf{55--60\% RH}, supporting 60\% as the practical
dampness threshold. Stratifying by meteorological sensitivity before prediction raised
F1 from 0.53 to \textbf{0.65 (bacteria) and 0.75 (mould)}. Culture-based counts, so the
non-viable fraction is invisible, and the facility-type classes are administrative rather
than physical. Directly relevant to anyone adding T/RH channels to an indoor node and
expecting them to carry bioaerosol information.}
{10.1155/ina/6724151}

\band{Exposure assessment \& modelling\hfill 16 records}{Amber}

\paperentry{Wang T et al.\ 2026}
{IAQMS-Industry: 100\,m PM$_{2.5}$ inside industrial parks}
{J Environ Sci (China)}
{Couples the regional NAQPMS to a city-scale chemical transport model with explicit
point-source locations and Gaussian plume dispersion, resolving PM$_{2.5}$ at
\textbf{100\,m} against a 3\,km regional baseline, evaluated at Beijing Yizhuang and
Tangshan. Normalised mean bias improves from $-$16.9\%\,--\,$-$7.7\% to
\textbf{3.1--6.2\%}, and industrial point sources account for \textbf{22.9--26.4\%} of
park PM$_{2.5}$ versus 1.6--13.7\% in the regional model. The mechanism is
straightforward --- coarse grids smooth the concentration gradient and over-disperse ---
but the size of the attribution error is the useful number. This is the exact error a
dense in-park sensor network would be deployed to measure, and the paper supplies the
modelled expectation to test against.}
{10.1016/j.jes.2025.12.020}

\paperentry{Boyle et al.\ 2026}
{Cancer-specific air-pollution exposome metrics for the continental US}
{Cancer Epidemiol Biomarkers Prev}
{Combines CACES land-use-regression criteria-pollutant surfaces with the RSEI
industrial-release model across US census tracts, 2005--2020, weighting each pollutant by
meta-analytic breast- and lung-cancer risk estimates to produce two composite metrics
(APBC, APLC). Both are strongly right-skewed --- some tracts an order of magnitude above
the mean --- peak in large cities and along transport corridors, differ from RSEI notably
in Los Angeles and New York, vary by race, ethnicity and housing, and cluster spatially.
No cancer outcome is analysed here: this is metric construction, and the meta-analytic
weights import the heterogeneity of their source studies into a tract-level surface that
then reads as precise. Useful as a template for building a risk-weighted composite from a
network, with that caveat attached.}
{10.1158/1055-9965.EPI-26-0493}

\paperentry{Arpornwichanop \& Chuersuwan 2026}
{Life-stage-resolved inhalation exposure to PM$_{2.5}$-bound elements, Greater Bangkok}
{Environ Geochem Health}
{Secondary-data screening at one fixed site (selected 24\,h filter days, Dec 2020--Sep
2021), combining seasonally weighted outdoor elemental concentrations with indoor
penetration, time-activity, respiratory-deposition and simulated-lung-fluid solubility
assumptions across pediatric life stages. Concentration and deposited-intake profiles
rank \textbf{Al, Fe, B, Zn} highest; the solubility- and toxicity-adjusted comparative
priority quotient ranks \textbf{As, Cd, Sb}. Modelled lifetime cancer risk spans
\textbf{1.11$\times$10$^{-6}$ to 2.78$\times$10$^{-5}$} across 0--100\% Cr(VI)
assumptions --- an order of magnitude driven entirely by a speciation parameter nobody
measured. The authors state explicitly that this is within-model prioritisation, not
measured child exposure. That honesty is unusual and worth noting.}
{10.1007/s10653-026-03455-7}

\paperentry{Xiao S et al.\ 2026}
{PAHs and oxygenated PAHs in Xi'an PM$_{2.5}$ under the clean air action}
{Environ Geochem Health}
{Winter PM$_{2.5}$-bound parent PAHs and OPAHs in Xi'an, a Fenwei Plain basin city, with
PMF source apportionment and toxic-equivalent concentrations. Wintertime PAHs are lower
than historical levels under the clean-air actions; biomass-burning-related mixed
combustion is the largest PMF factor. During haze episodes both classes rise, but the
increase is \textbf{statistically significant for OPAHs and not for parent PAHs}, and
\textbf{OPAHs --- BcdPQ in particular --- dominate the toxic-equivalent burden}. Regional
pollution is in-basin accumulation plus short-range transport from northern Shaanxi
industry, modulated by synoptic pressure pattern. The finding that matters: a control
programme that lowers parent PAHs need not lower the oxygenated derivatives that carry
the toxicity, and routine monitoring measures the former.}
{10.1007/s10653-026-03449-5}

\paperentry{G\"urler Aky\"uz \& Aky\"uz 2026}
{Size-resolved elemental composition in a coastal urban-industrial region}
{J Environ Sci Health A}
{Seasonal PM$_{2.5}$ and PM$_{2.5-10}$ elemental composition at Zonguldak, T\"urkiye,
where coal mining, coal combustion, iron-steel with coke production, maritime fuel oil,
traffic and marine/crustal sources overlap. Heating-period PM$_{2.5}$ carries elevated
\textbf{As, Sb, Cd, Pb, Se, V, Ni, Cr}; non-heating PM$_{2.5-10}$ is crustal, marine and
resuspended dust with a persistent industrial residue. Enrichment factors against Ti show
strong anthropogenic enrichment for combustion and metallurgy and near-natural levels for
crustal and marine elements; PCA separates coal combustion, iron-steel, fuel-oil/shipping,
non-exhaust traffic, crustal dust, marine aerosol and coal fly ash, and discriminant
analysis confirms the season/size separation. Conventional in method, valuable as a
composition baseline for a source mixture that is rarely characterised together.}
{10.1080/10934529.2026.2724190}

\paperentry{Guo \& Sudo 2026}
{Global online PAH simulation in the CHASER chemistry-climate model}
{Atmosphere}
{First global online simulation of naphthalene, phenanthrene, pyrene and benzo[a]pyrene in
CHASER V4.0, coupling soot adsorption $K_{SA}(T)$ with organic-matter absorption
$K_{OA}(T)$ under an extended Dachs--Eisenreich dual-mode scheme. Evaluated at urban,
background and Arctic stations: pooled spatial $R$ \textbf{0.78} on a log$_{10}$ scale,
\textbf{FAC2 80\%} for BaP excluding sub-grid hotspots, seasonal $R$ 0.62--0.93. The
$K_{SA}$ term fixes the long-standing under-prediction of particle-bound fraction at cold
high latitudes, raising polar $\phi_{\mathrm{BaP}}$ by up to $\sim$0.6, and elevated
elemental carbon extends BaP lifetime to \textbf{1.72\,d} against 0.57--0.72\,d for
lighter congeners, sustaining transoceanic outflow. The dependence on modelled EC is also
the weak point: the partitioning result is only as good as the soot field underneath it.}
{10.3390/atmos17090861}

\paperentry{Klebach et al.\ 2026}
{MSA and sulfuric acid in tropical Indo-Pacific aerosol from aircraft}
{Atmos Chem Phys}
{CAFE-Pacific, HALO aircraft, January--February 2024, nitrate chemical-ionisation mass
spectrometry of methanesulfonic and sulfuric acid. Marine boundary layer concentrations
reach \textbf{4$\times$10$^7$\,cm$^{-3}$ MSA and 6$\times$10$^7$\,cm$^{-3}$ SA}; the
MSA/SA ratio increases with altitude in the lower free troposphere, consistent with
temperature-dependent DMS oxidation. Two flights show marine deep convection lofting DMS
to the upper troposphere with particle formation and growth after \textbf{10--20\,h of OH
exposure}, matching the kinetically modelled DMS lifetime. MSA frequently exceeds SA,
which the current generation of new-particle-formation schemes does not anticipate.
Aircraft inlets heat and evaporate acidic particles, so the gas/particle split at altitude
is instrumental as much as atmospheric --- the authors are explicit about this.}
{10.5194/acp-26-12355-2026}

\paperentry{Papetta et al.\ 2026}
{Source-dependent optical and mineral signatures of Mediterranean dust outbreaks}
{Atmos Chem Phys}
{Four major 2021--2022 dust outbreaks from the Eastern, Western and Central Sahara and the
Middle East, tracked across \textbf{24 AERONET stations} with IASI and MODIS MIDAS dust
optical depth and HYSPLIT back-trajectories. Saharan events were coarse and scattering;
the Middle East event carried \textbf{finer, more absorbing particles}, likely
anthropogenically influenced. DOD-to-AOD ratio separates them --- only one East-Central
Saharan event stayed above 0.8 (near-pure dust), while the Middle East event had the
lowest ratios, indicating mixing with anthropogenic or marine aerosol. A Cyprus in-situ
case study with elemental and absorption measurements anchors the remote-sensing
inference. For network work the transferable idea is the DOD/AOD purity screen: it is a
cheap way to decide whether a coarse-mode excursion is genuinely dust before attributing
it.}
{10.5194/acp-26-12395-2026}

\paperentry{Liu S et al.\ 2026}
{Terrain blocking and monsoon coupling in North China Plain PM$_{2.5}$ and O$_3$}
{J Environ Sci (China)}
{Quantifies mountain-blocking effects on pollutant dispersion in the semi-closed North
China Plain now that anthropogenic emissions have fallen far enough for terrain and
meteorology to dominate the residual pattern. Within 120\,km buffers of the major ranges,
concentrations rise with altitude and then fall sharply above a threshold; the
\textbf{effective blocking altitude is 22--929\,m for winter PM$_{2.5}$ and 19--1060\,m
for summer O$_3$}. Concentrations rank plains $>$ platforms $>$ hills $>$ mountains, and
winter PM$_{2.5}$ and summer O$_3$ high-value zones aggregate in \emph{opposite}
north--south directions. Descriptive rather than causal --- emissions and terrain
co-vary strongly in the NCP --- but the altitude thresholds are a concrete siting
constraint for any regional network.}
{10.1016/j.jes.2026.01.013}

\paperentry{Stergiou et al.\ 2026}
{Spatiotemporal graph neural network bias correction for CMAQ over CONUS}
{Atmos Pollut Res}
{Ships on metadata only. Crossref carries no abstract for this DOI and the record is not
in PubMed, Europe PMC or Semantic Scholar as of this run, so nothing beyond title,
authors and venue can be attributed to it. Counted and listed; no numbers claimed.
Flagged for the abstract retry list --- graph-structured bias correction of a regulatory
CTM against monitor networks is directly on this watch's topic and should be summarised
once the abstract deposits.}
{10.1016/j.apr.2026.103191}

\paperentry{Xu L et al.\ 2026}
{Marine aging compacts soot and thickens its coating}
{J Environ Sci (China)}
{TEM single-particle analysis of soot collected over the Bohai Sea and the Northern and
Southern Yellow Sea after outflow from East Asia. \textbf{Over 98\%} of soot-containing
particles are coated, mean mixing-state index \textbf{0.83}. Fractal dimension rises along
the transport path --- 1.84 $\pm$ 0.05 (N. Yellow Sea), 1.90 $\pm$ 0.08 (Bohai), 1.96
$\pm$ 0.07 (S. Yellow Sea) --- indicating progressive structural compaction, while
$D_p/D_{core}$ of 3.9--5.3 exceeds the 3.54 typical of continental regional transport.
Compact, uniformly coated, thickly shelled soot absorbs substantially more than the
lacy fresh particle that optical models assume. Any absorption-based black-carbon
instrument calibrated on fresh soot will misreport aged marine air, and this paper
quantifies by how much the morphology has moved.}
{10.1016/j.jes.2026.01.047}

\paperentry{Zhang N et al.\ 2026}
{Aging, composition and light absorption of nitrogen-containing organics in PM$_1$}
{J Environ Sci (China)}
{59 Shanghai PM$_1$ samples analysed by UHPLC-Orbitrap MS. Nitrogen-containing organics
are \textbf{31\% of detected species in ESI$-$ and 64\% in ESI$+$}. Aging reduces
molecular diversity: CHON$-$ compounds fall in both number and mass fraction with aging
degree, and carboxylic-rich alicyclic molecules drop from \textbf{37.7\% to 21.2\%} of the
CHON$-$ contribution. Mass absorption efficiency at 365\,nm declines with aging alongside
the Bep/(Bep+Bap) ratio and relative humidity, pointing to aqueous-phase oxidation as the
dominant mechanism, with CHON$-$ and CHN$+$ identified as the principal chromophores.
Semi-quantitative by design --- ESI response factors are compound-specific and
uncalibrated --- so the fractions are ranks, not concentrations.}
{10.1016/j.jes.2025.12.003}

\paperentry{Xu X et al.\ 2026}
{Ammonia enhances isoprene SOA through gas-phase cluster formation}
{J Environ Sci (China)}
{Smog-chamber series on isoprene photooxidation by OH with and without NH$_3$, with
nitrate-CIMS for the gas phase and HR-TOF-AMS for the particle phase. NH$_3$ increased
SOA mass and accelerated isoprene oxidation. The mechanism is not the expected
particle-phase acid--base reaction: CIMS shows NH$_3$ reacting homogeneously in the
\emph{gas} phase with organic acids and low-volatility oxygenated organic molecules to
form extremely and ultra-low-volatility \textbf{NH$_3$--OOM clusters}, which then nucleate
and condense. Quantum chemistry confirms spontaneous hydrogen-bonded association with
RCOOH, R-OOH and R-OH groups. \textbf{Cluster formation accounts for 78\% of the SOA
enhancement.} Chamber conditions with elevated NH$_3$; the relevance is to agricultural
and livestock regions where NH$_3$ and biogenic VOC co-occur.}
{10.1016/j.jes.2026.01.041}

\paperentry{Kong J et al.\ 2026}
{Triplet-state brown carbon as an endogenous oxidant}
{J Environ Sci (China)}
{Photochemistry of 1,4-naphthoquinone as a brown-carbon proxy in deoxygenated aqueous
systems under 355\,nm irradiation, resolving two long-lived $^3(n,\pi^*)$ triplets
($^3$NQ* and $^3$OH-NQ*). $^3$NQ* abstracts hydrogen from organic substrates
\textbf{without any external oxidant}, with a rate hierarchy set by bond dissociation
energy and one-electron oxidation potential: alkanes $\sim$10$^4$, oxygenated alkanes
$\sim$10$^5$, phenols $\sim$10$^6$\,M$^{-1}$\,s$^{-1}$. Self-photoconversion regenerates
a second triplet, amplifying the oxidative cascade. The claim is that
$^3$BrC* is the dominant phenol sink in BrC-rich haze under acidic, O$_2$-limited
conditions --- which is exactly the condition inside a haze particle and exactly the
condition current models do not represent. Proxy compound, deoxygenated system; the
rate constants are the transferable output.}
{10.1016/j.jes.2025.10.037}

\paperentry{Wang J et al.\ 2026}
{Water-soluble aerosol matter photosensitises guaiacol into aqueous SOA}
{J Environ Sci (China)}
{Mechanism and kinetics of aqSOA formation from guaiacol mediated by photosensitive
water-soluble inorganic and organic matter extracted from real aerosol. Photolysis of the
extract generates both reactive oxygen species and triplet WSOM; \textbf{$^3$WSOM* is the
predominant initiator} of guaiacol transformation, with steady-state concentration
\textbf{5.35 $\pm$ 0.12 $\times$10$^{-13}$\,mol/L}, generation rate
3.43$\times$10$^{-7}$\,mol\,L$^{-1}$\,s$^{-1}$, quantum yield \textbf{4.86\%} and a rate
constant with guaiacol above 10$^9$\,L\,mol$^{-1}$\,s$^{-1}$. ROS then convert soluble
intermediates into the insoluble ethers and low-volatility phenols that constitute the
aqSOA. Bulk aqueous extracts are not cloud or particle water --- ionic strength differs by
orders of magnitude --- which is the standing limitation of this whole method.}
{10.1016/j.jes.2025.11.050}

\paperentry{Solaiman et al.\ 2026}
{Nighttime NO$_3$ oxidation of myrtenal: SOA yields and HOMs}
{Environ Sci Atmos}
{Chamber study of secondary organic aerosol from the gas-phase oxidation of myrtenal, a
biogenic aldehyde, by NO$_3$ radicals, reporting SOA yields, chemical composition and
highly oxygenated molecules. The Crossref deposit truncates the abstract mid-sentence and
no full record was retrievable from Semantic Scholar on this run, so the yields and HOM
distributions cannot be quoted. Counted and listed; the nighttime NO$_3$ channel for
oxygenated biogenics is under-represented in this watch's corpus and the record is worth
revisiting once the full abstract deposits.}
{10.1039/d6ea00092d}

\band{Burden, policy \& mitigation\hfill 3 records}{Coral}

\paperentry{Zhao N et al.\ 2026}
{Attainment pressure under China's revised GB 3095-2026 standard in Shandong}
{J Hazard Mater}
{Daily observations of six criteria pollutants across 16 Shandong cities, 2015--2025,
condensed into a relative attainment pressure index combining concentration gap and
additional non-attainment days, with random forest, SHAP and monotonicity-constrained
XGBoost on socioeconomic predictors. \textbf{PM$_{2.5}$, PM$_{10}$, SO$_2$, NO$_2$ and CO
fell 41.0--79.6\%; O$_3$ did not move, staying at 166--191\,$\mu$g\,m$^{-3}$.} Under the
transition-stage limits a western high-pressure belt emerges --- Heze 98.7, Dezhou 91.3,
Liaocheng 89.1. Civilian vehicle ownership and secondary-industry share are the leading
predictors for particulates and NO$_2$; the socioeconomic set explains PM$_{2.5}$ far
better than O$_3$. Even under the deep-transition pathway \textbf{eight of sixteen cities
remain above 30\,$\mu$g\,m$^{-3}$ in 2030}. One province, and the scenario arithmetic
assumes the historical predictor--concentration relation holds forward.}
{10.1016/j.jhazmat.2026.143460}

\paperentry{Li Y et al.\ 2026}
{Multidimensional inequality of CO$_2$ and air-pollutant footprints in China}
{Environ Sci Technol}
{Consumption-based per-capita footprint accounting for CO$_2$ and air pollutants across
2020--2060 under the dual-carbon policy. The top 10\% carry \textbf{28.5\,t CO$_2$ and
50.6\,kg NO$_x$} per capita against \textbf{1.2\,t and 4.7\,kg} for the bottom 50\%. The
structural finding is that \textbf{over 90\% of footprint inequality is intraprovincial},
not between provinces, and the gap widens to 2060 --- which means province-level policy
instruments are aimed at the wrong axis. Composition differs by stratum: the top 1\%
footprint is investment-driven (96\% of SO$_2$), the bottom 50\% consumption-driven
(88\% of NH$_3$). Input--output footprint accounting carries well-known allocation
sensitivity, and the 2060 horizon is a scenario, not a projection.}
{10.1021/acs.est.6c04314}

\paperentry{Shi et al.\ 2026}
{Carbon--air pollution cobenefits and emission leakage in China's steel industry}
{Environ Sci Technol}
{Unit-level emission database covering \textbf{2{,}866 enterprises and $\sim$4{,}030
units}, relating regional energy-industrial configuration to mitigation priorities. Steel
CO$_2$ peaked at 1.99\,Gt in 2019 while SO$_2$ and PM fell \textbf{more than 80\%} over
2011--2022 --- the decoupling this watch keeps seeing. Over 2023--2050 the priorities
diverge: production restraint gives the largest cumulative CO$_2$ reduction (22.1\%)
while \textbf{obsolete-capacity retirement dominates air-pollutant reduction (22.3\%)},
so a carbon-optimal pathway is not a PM-optimal one. Regionally differentiated pathways
induce \textbf{spatial emission leakage}, stronger for CO$_2$ (17.39\% in 2030, falling to
2.38\% by 2050) than for pollutants, and national mean pollutant concentrations fall only
5.4--12.3\% by 2030 relative to 2022. Sector-specific and inventory-based.}
{10.1021/acs.est.6c08461}

\clearpage

% ==================================================================== TABLE
\section*{\sffamily\bfseries\color{Deep}Corpus register}
\vspace{-1.5mm}{\color{Amber}\rule{\linewidth}{1.1pt}}\vspace{2.5mm}

{\fontsize{6.7}{8.4}\selectfont
\setlength{\tabcolsep}{2.4pt}
\begin{longtable}{@{}>{\RaggedRight}p{25mm}>{\RaggedRight}p{41mm}>{\RaggedRight}p{27mm}%
>{\RaggedRight}p{23mm}>{\RaggedRight}p{18mm}>{\RaggedRight}p{15mm}@{}}
\toprule
\textsf{\textbf{First author}} & \textsf{\textbf{Focus}} & \textsf{\textbf{Journal}} &
\textsf{\textbf{Design}} & \textsf{\textbf{Metric}} & \textsf{\textbf{Region}}\\
\midrule
\endfirsthead
\toprule
\textsf{\textbf{First author}} & \textsf{\textbf{Focus}} & \textsf{\textbf{Journal}} &
\textsf{\textbf{Design}} & \textsf{\textbf{Metric}} & \textsf{\textbf{Region}}\\
\midrule
\endhead
\multicolumn{6}{@{}l@{}}{\cellcolor{Mist}\textsf{\textbf{\textcolor{Deep}{Neuro / mental health}}}}\\[0.6mm]
\href{https://doi.org/10.1001/jamanetworkopen.2026.25015}{Dales} & Wildfire-sourced PM$_{2.5}$ and ED visits for migraine and other primary headache & \textit{JAMA Netw Open} & Case-crossover & PM$_{2.5}$ (wildfire) & Canada\\
\href{https://doi.org/10.1093/ageing/afag215}{Glover} & Umbrella meta-review of urban environmental exposures and cognitive health & \textit{Age Ageing} & Umbrella review / meta-analysis & PM$_{2.5}$ & Global\\
\href{https://doi.org/10.1016/j.nsa.2026.107019}{Goossens} & In utero air pollution exposure and pre-/postnatal brain development & \textit{Neurosci Appl} & Narrative review & PM$_{2.5}$ / PM$_{10}$ / PM$_{\mathrm{coarse}}$ & Global\\
\href{https://doi.org/10.1016/j.isci.2026.116617}{Ma} & PM$_{2.5}$-induced neurotoxicity and depression: exposure risks and mechanisms & \textit{iScience} & Narrative review & PM$_{2.5}$ + components & Global\\
\href{https://doi.org/10.3760/cma.j.cn112150-20250804-00752}{Pan} & Maternal PM exposure and offspring neurodevelopment: effects and mechanisms & \textit{Zhonghua Yu Fang Yi Xue Za Zhi} & Narrative review & PM (bulk) & China\\
\addlinespace[0.9mm]
\multicolumn{6}{@{}l@{}}{\cellcolor{Mist}\textsf{\textbf{\textcolor{Teal}{Cardiovascular \& metabolic}}}}\\[0.6mm]
\href{https://doi.org/10.3760/cma.j.cn112338-20251120-00843}{Huang} & PM$_{2.5}$ and dyslipidaemia in rural Chinese children; dietary effect modification & \textit{Zhonghua Liu Xing Bing Xue Za Zhi} & Cross-sectional (multistage) & PM$_{2.5}$ & China\\
\href{https://doi.org/10.32394/pe/218966}{Kazmierska} & Air pollution and cardiovascular disease: review of current data & \textit{Przegl Epidemiol} & Narrative review & PM$_{2.5}$ / PM$_{10}$ & Global\\
\href{https://doi.org/10.1007/s00125-026-06814-2}{Xi} & Air pollution, built environment and temperature vs. prediabetes phenotypes & \textit{Diabetologia} & Prospective cohort & PM$_{2.5}$ & Germany\\
\href{https://doi.org/10.1161/JAHA.125.046917}{Yang} & Short-term PM$_{2.5}$ and hospitalisation/ED use in atrial fibrillation & \textit{J Am Heart Assoc} & Case-crossover & PM$_{2.5}$ (wildfire-influenced) & USA\\
\addlinespace[0.9mm]
\multicolumn{6}{@{}l@{}}{\cellcolor{Mist}\textsf{\textbf{\textcolor{Sky}{Reproductive \& developmental}}}}\\[0.6mm]
\href{https://doi.org/10.1016/j.jes.2025.10.043}{Chen} & Prenatal air pollution, fetal growth trajectory phenotypes and child neurodevelopment & \textit{J Environ Sci (China)} & Prospective cohort & PM$_{2.5}$ / PM$_{10}$ & China\\
\href{https://doi.org/10.1016/j.ecoenv.2026.120566}{Huang} & Joint prenatal bisphenol + air-pollutant mixture exposure and preschool IQ & \textit{Ecotoxicol Environ Saf} & Prospective cohort (mixtures) & PM$_{2.5}$ / PM$_{10}$ & China\\
\href{https://doi.org/10.1007/s40572-026-00552-8}{Peterson} & Ambient air pollution, PAHs and pubertal development: critical review & \textit{Curr Environ Health Rep} & Critical review & PM$_{2.5}$ components & Global\\
\href{https://doi.org/10.1016/j.envres.2026.125313}{Xie} & Long-term PM$_{2.5}$ and timing of pubertal development in boys and girls & \textit{Environ Res} & Prospective cohort & PM$_{2.5}$ & China\\
\addlinespace[0.9mm]
\multicolumn{6}{@{}l@{}}{\cellcolor{Mist}\textsf{\textbf{\textcolor{Sage}{Respiratory \& allergic}}}}\\[0.6mm]
\href{https://doi.org/10.1007/s00484-026-03289-z}{Bayar Muluk} & Climate change and otolaryngology: airway disease impacts and emission reduction & \textit{Int J Biometeorol} & Narrative review & PM$_{2.5}$ & Global\\
\href{https://doi.org/10.4103/mgr.MEDGASRES-D-25-00164}{Liao} & Short-term air pollution, respiratory symptoms and lung function in asthma & \textit{Med Gas Res} & Panel study & PM$_{2.5}$ / PM$_{10}$ & China\\
\href{https://doi.org/10.3390/epidemiologia7040097}{Sanchez de Prada} & Ambient PM and asthma-coded admissions in children 0-3 y, Spain & \textit{Epidemiologia} & Ecological panel & PM$_{2.5}$ / PM$_{10}$ & Spain\\
\href{https://doi.org/10.1016/j.jhazmat.2026.143076}{Yang} & Allergenic vs. toxic risk of airborne seed hairs compared with ambient PM$_{2.5}$ & \textit{J Hazard Mater} & Field measurement & PM$_{2.5}$ + bioaerosol & China\\
\addlinespace[0.9mm]
\multicolumn{6}{@{}l@{}}{\cellcolor{Mist}\textsf{\textbf{\textcolor{Violet}{Mechanistic toxicology}}}}\\[0.6mm]
\href{https://doi.org/10.1016/j.jes.2025.12.074}{Ding} & PM$_{2.5}$ promotes vascular calcification via miR-27b-5p/HES1 signalling & \textit{J Environ Sci (China)} & In vitro + organoid + murine & PM$_{2.5}$ & China\\
\href{https://doi.org/10.1007/s00360-026-01698-5}{Maraschi} & Osmoregulatory disruption and metal bioaccumulation from settleable atmospheric PM & \textit{J Comp Physiol B} & Controlled exposure (ecotox) & Settleable APM & Brazil\\
\href{https://doi.org/10.1016/j.jes.2025.11.018}{Wu} & Macrophage NLRP3 inflammasome in real-ambient PM-induced hepatic glucose dysregulation & \textit{J Environ Sci (China)} & Real-ambient murine + co-culture & Real-ambient PM & China\\
\addlinespace[0.9mm]
\multicolumn{6}{@{}l@{}}{\cellcolor{Mist}\textsf{\textbf{\textcolor{Clay}{Occupational \& indoor}}}}\\[0.6mm]
\href{https://doi.org/10.3389/fmedt.2026.1796828}{Chen} & Nested suction instrument to cut intraoperative surgical-smoke PM$_{2.5}$ & \textit{Front Med Technol} & Bench / device experiment & PM$_{2.5}$ (surgical smoke) & China\\
\href{https://doi.org/10.1093/annweh/wxag055}{Cole} & Determinants of respirable crystalline silica exposure in Australian tunnelling & \textit{Ann Work Expo Health} & Observational exposure study & Respirable dust / RCS & Australia\\
\href{https://doi.org/10.1016/j.jhazmat.2026.143081}{Ding} & Stage-resolved hair Zn-isotope dynamics as a firefighting combustion-exposure tracer & \textit{J Hazard Mater} & Repeated-measures biomarker & Combustion-derived PM & China\\
\href{https://doi.org/10.3390/jox16040119}{Ramsden} & Xenobiotic hazards in aircraft cabin air & \textit{J Xenobiot} & Review + risk economics & Ultrafine particles & Global\\
\addlinespace[0.9mm]
\multicolumn{6}{@{}l@{}}{\cellcolor{Mist}\textsf{\textbf{\textcolor{Amber}{Sensing, forecasting \& instrumentation}}}}\\[0.6mm]
\href{https://doi.org/10.1016/j.jenvman.2026.130515}{Shi} & Adaptive spatiotemporal graph transformer for multi-site multi-step PM$_{2.5}$ forecasting & \textit{J Environ Manage} & Deep-learning model & PM$_{2.5}$ & China / India\\
\href{https://doi.org/10.3390/bios16070380}{Tsuchiya} & Seasonal and indoor-outdoor characteristics of fluorescent aerosol particles in offices & \textit{Biosensors} & Real-time sensor campaign & Bioaerosol / coarse & Japan\\
\addlinespace[0.9mm]
\multicolumn{6}{@{}l@{}}{\cellcolor{Mist}\textsf{\textbf{\textcolor{Amber}{Exposure assessment \& modelling}}}}\\[0.6mm]
\href{https://doi.org/10.1029/2025GH001621}{Solorzano-Araque} & Multiscale exposure modelling: LOTOS-EUROS + mobility surveys + particle toxicity & \textit{GeoHealth} & CTM + mobility integration & PM$_{2.5}$ / PM$_{10}$ & Colombia\\
\addlinespace[0.9mm]
\multicolumn{6}{@{}l@{}}{\cellcolor{Mist}\textsf{\textbf{\textcolor{Coral}{Burden, policy \& mitigation}}}}\\[0.6mm]
\href{https://doi.org/10.1002/hsr2.72886}{Karamzad} & Ischaemic heart disease burden in MENA 1990-2021; PM attributable fraction 28.2\% & \textit{Health Sci Rep} & GBD secondary analysis & Ambient PM & MENA\\
\href{https://doi.org/10.1371/journal.pone.0347989}{Rees} & HEPA filtration to reduce respiratory infection in care homes & \textit{PLoS One} & RCT process evaluation & Filtered PM & UK\\
\href{https://doi.org/10.1016/j.jes.2026.01.009}{Shan} & Policy-driven emission reductions and post-pandemic rebound, high-resolution inventories & \textit{J Environ Sci (China)} & Emission inventory & PM emissions & China\\
\addlinespace[0.9mm]
\multicolumn{6}{@{}l@{}}{\cellcolor{Mist}\textsf{\textbf{\textcolor{Slate}{Other clinical endpoints}}}}\\[0.6mm]
\href{https://doi.org/10.3390/dj14070440}{Lupita} & PM$_{2.5}$, preventive-care capacity and caries/edentulism burden across the EU & \textit{Dent J} & Ecological panel & PM$_{2.5}$ & EU-27\\
\href{https://doi.org/10.1016/j.jdin.2026.06.008}{Park} & Long-term PM exposure and risk of androgenetic alopecia & \textit{JAAD Int} & Nationwide cohort & PM$_{2.5}$ / PM$_{10}$ & South Korea\\
\href{https://doi.org/10.3171/2026.3.PEDS2615}{Sun} & Environmental exposure and area deprivation vs. paediatric hydrocephalus outcomes & \textit{J Neurosurg Pediatr} & Registry cohort + geospatial & PM$_{2.5}$ (satellite-derived) & USA\\
\addlinespace[0.9mm]
\bottomrule
\end{longtable}}


\vspace{4mm}
\begin{keybox}
{\sffamily\bfseries\footnotesize\color{Deep}Provenance and method}\\[1.2mm]
{\fontsize{7.2}{8.8}\selectfont
Entry window \textbf{1 September 2026}, single day, opened against a
\texttt{last\_entry\_date} of 31 August so no day is skipped. Sources queried on
3 September 2026 (local 20:5x CDT, so under this pipeline's 22:00 rule the current day is
not yet buildable and 1--2 September are the owed issues).
\textbf{PubMed} E-utilities on \texttt{[EDAT]} 2026/09/01: 26 hits on the health axis and
5 on the instrumentation axis, run as separate queries.
\textbf{Europe PMC} \texttt{CREATION\_DATE}: 11 hits.
\textbf{Crossref} by-ISSN, windowed on \texttt{created}: 97 deposits across the tracked
atmospheric and exposure journals --- the channel that supplies \textit{Atmos.\ Chem.\
Phys.}, \textit{Atmosphere} and \textit{Environ.\ Sci.\ Atmos.}, none of which PubMed
indexes.
\textbf{Consensus} on the sensor-calibration axis: 3 hits, \textbf{all 3 already in the
store}.
\textbf{OpenAlex} returned HTTP 403 (plan-gated date filter) and \textbf{arXiv} HTTP 429
on repeated attempts; both legs are recorded as failed, not as empty.
After collapsing cross-source duplicates by DOI and normalised title, 135 unique records
were considered; \textbf{4 were already shipped} (all on 31 August, caught on the PubMed
identifier leg), 33 were screened in scope and 98 were screened out as off-topic ---
predominantly aquatic and soil chemistry, microbiology, meteorology and non-particulate
building-science deposits carried in by the by-ISSN sweep.
Subtopic, design, particle-metric and geography labels are single-label by dominant
contribution and are a judgement, not a controlled vocabulary. Effect estimates are
transcribed verbatim from abstracts. Two records (\texttt{10.1016/j.apr.2026.103191},
\texttt{10.1039/d6ea00092d}) ship on metadata only and no numbers are attributed to them.
All DOI links resolve to the publisher of record; abstracts remain the copyright of their
respective publishers.\par}
\end{keybox}

\end{document}
